\subsection{Polarizations, ordinary replacements, and descent}
\label{prog:foundations}

The program constructions below use the analytic cone
\(\overline{\mathrm{NA}}(T)\), dual to the nef cone in
\(H^{1,1}_{\BC}(T,\R)\) for spaces with rational singularities
\cite[Proposition~2.4]{DHP}. Degrees of global line bundles form a
possibly smaller numerical space. We pass from the analytic cone to
projective geometry by producing an actual relatively ample line.
The ensuing descent statements apply on singular working models as
well as on smooth initial spaces.

\begin{lemma}[Degree-zero Bott--Chern descent]
\label{prog:bc-descent}
Let \(f:T\to Z\) be a proper surjective morphism with connected fibers
between normal compact complex spaces with rational singularities.
Suppose either
\begin{enumerate}
\item \(f\) is bimeromorphic and both spaces belong to Fujiki's class
\(\mathcal C\); or
\item \(f\) is projective and an effective rational boundary
\(\Delta\) makes \((T,\Delta)\) klt with
\(-(K_T+\Delta)\) relatively nef and big.
\end{enumerate}
Pullback is injective on both \(H^2(-,\R)\) and
\(H^{1,1}_{\BC}(-,\R)\). In either group its image consists exactly
of the classes with degree zero on every curve contracted by \(f\).
\end{lemma}

\begin{proof}
These are, respectively, the two cases of
\cite[Lemma~2.6(1),(2)]{DH2024}. The first case is birational descent
between spaces with rational singularities and does not require a
projectivity criterion. The second case uses relative vanishing for
the ordinary rational klt pair and the projective relative curve
space. In particular, it can be applied to a projective log-Fano
contraction before constructing generalized data on its base.
\end{proof}

In applications to an ordinary birational negative contraction,
the rationality hypotheses are automatic. The source is locally klt.
For a dlt pair, first decrease its rational Cartier floor slightly;
relative antiampleness persists. Relative vanishing gives
\(R^if_*\OO_T=0\) for \(i>0\)
\cite[Theorem~5.2]{FujinoAnalyticBCHM}. Composing a resolution of \(T\)
with \(f\), the Leray spectral sequence proves that \(Z\) has rational
singularities. The same reasoning works when the klt boundary is
chosen separately on finitely many Stein base neighborhoods. For
fiber-type contractions in case~(ii), rationality of the base follows
from \cite[Lemma~8.8(i)]{DHP}; thus that application also precedes
the construction of generalized data on the base.

\begin{lemma}[Projectivity from a rational line detecting the ray]
\label{prog:detected-projectivity}
Let $f:T\to Z$ be a proper contraction between normal compact
K\"ahler spaces, with $T$ having rational singularities. Suppose
$\gamma$ is K\"ahler on $Z$ and
\[
 \alpha=f^*\gamma,\qquad
 \overline{\mathrm{NA}}(T)\cap\alpha^\perp=R
\]
is a nonzero ray. If a global rational line bundle $D$ has
$D\cdot R<0$, then $-D$ is relatively ample and $f$ is projective.
\end{lemma}
\begin{proof}
Fix a K\"ahler class $\omega$ on $T$ and normalize $\overline{\mathrm{NA}}(T)$ by
$\omega\cdot z=1$. The resulting slice $S$ is compact. Its unique
point on $R$ has negative $D$-degree, so $d=c_1(D)$ is negative on
a neighborhood of that point in $S$. On the complementary compact
set, $\alpha$ has a strictly positive minimum, while $d$ is bounded.
For all sufficiently small $s>0$, the class $\alpha-sd$ is therefore
strictly positive on $S$ and is K\"ahler by cone duality.

Choose an integer $m>0$ making $mD$ integral. Then
\[
 c_1(-mD)+f^*((m/s)\gamma)=(m/s)(\alpha-sd)
\]
is K\"ahler. On a neighborhood in the base with a potential for
$\gamma$, absorb that potential in a smooth metric on the same global
line bundle $-mD$. Its curvature is positive on the fibres. The analytic
positive-line-bundle criterion and the fibrewise relative-ampleness
criterion make $-mD$ relatively ample. The line bundle was global
throughout, so no gluing of unrelated local polarizations is required.
\end{proof}


An analytic contraction supplied by a supporting class therefore becomes
projective as soon as one rational line has negative degree on its ray.
For a small contraction, that same line will define the flip globally.
The local ordinary adjoints used to prove finite generation need not
glue; their section algebras will be Veronese subalgebras of one algebra.

\begin{lemma}[Divisor representatives over a Stein base]
\label{prog:stein-representatives}
Let $\pi:Y\to U$ be a projective surjective morphism with connected fibres
between normal irreducible complex spaces, where $U$ is Stein, and let
$\mathcal L$ be a holomorphic line bundle on $Y$.
If $\pi$ is bimeromorphic, $\mathcal L$ has a holomorphic section
which is not identically zero, and hence an effective Cartier divisor
representative. More generally, after shrinking $U$ around any specified
point, $\mathcal L$ has a nonzero meromorphic section.
\end{lemma}

\begin{proof}
Suppose first that $\pi$ is bimeromorphic. Grauert's proper direct-image
theorem makes $\mathcal F=\pi_*\mathcal L$ coherent. On the nonempty
open subset $U^\circ$ where $\pi$ is an isomorphism, it is a line bundle.
Choose $u_*\in U^\circ$. Cartan's theorem~A says that global sections
generate $\mathcal F_{u_*}$, so one global section has nonzero image in
$\mathcal F_{u_*}/\mathfrak m_{u_*}\mathcal F_{u_*}$.
Under
\[
 H^0(U,\pi_*\mathcal L)=H^0(Y,\mathcal L)
\]
it gives a section $s$ nonzero at the unique point above $u_*$.
Its zero divisor is effective Cartier and represents $\mathcal L$.
The section may vanish elsewhere; only its meromorphic inverse is
being used to obtain a meromorphic trivialization.

For the general case, fix $u\in U$ and $y\in\pi^{-1}(u)$, and choose
a $\pi$-ample line bundle $\mathcal H$. Relative Serre generation,
after shrinking to a Stein neighbourhood of $u$, gives an integer
$n$ for which both $\mathcal L\otimes\mathcal H^n$ and
$\mathcal H^n$ are relatively generated. Their proper direct images
are coherent. Cartan's theorem~A and the evaluation maps supply
sections $s$ and $t$ of these two bundles which are both nonzero at
$y$. The quotient $s/t$ is the required nonzero meromorphic section
of $\mathcal L$. At no point is $Y$ assumed Stein.
\end{proof}

\begin{proposition}[A global algebra for a detected flip]
\label{prog:detected-flip}
Let $(T,B+\mathbf M)$ be a gklt pair on a normal compact K\"ahler
space which is globally strongly $\Q$-factorial, and let $A\in H^{1,1}_{\BC}(T,\R)$
be its adjoint class. Let $f:T\to Z$ be a projective small contraction
onto a normal compact K\"ahler space. Suppose that the classes of all
$f$-vertical curves lie on one ray $R\subset\overline{\mathrm{NA}}(T)$, that $A\cdot R<0$,
and that a global line bundle $L$ satisfies $L\cdot R<0$.

Then $\bigoplus_{m\geq0}f_*L^m$ is locally finitely generated. The
relative Proj of a sufficiently divisible Veronese is a projective
small contraction $f^+:T^+\to Z$, where $T^+$ is compact K\"ahler and
globally strongly $\Q$-factorial. The reflexive transform $L^+$ is a
rational line bundle, with a positive power equal to the relatively
ample tautological bundle. The transformed boundary and the same
b-nef datum define a gklt pair on $T^+$, with relatively K\"ahler
adjoint $A^+$. This is the generalized flip, with the usual
strict discrepancy increase at centres in either exceptional locus.

Moreover, if $t=(A\cdot C)/(L\cdot C)>0$ for an $f$-vertical curve $C$,
then there is a unique $\eta\in H^{1,1}_{\BC}(Z,\R)$ such that
\[
 A=t\,c_1(L)+f^*\eta,\qquad
 A^+=t\,c_1(L^+)+(f^+)^*\eta.
\]
Every class on $T$ has the corresponding forward transport with an
actual exceptional-divisor correction. No assertion of surjectivity
on Bott--Chern groups is required.
\end{proposition}

\begin{proof}
Since $f$ has a global relatively ample bundle and all its vertical
curves lie on $R$, numerical invariance of ampleness on the projective
fibres shows that $-L$ is $f$-ample.

\smallskip\noindent
\emph{A local rational ordinary adjoint.}
Fix $z\in Z$ and a relatively compact Stein neighbourhood $U$.
Take a projective log resolution $p:V\to T_U$ carrying $\mathbf M$,
and put $\rho=f p$. Write its actual structure boundary as $B_V$, so
\[
 [K_V+B_V]+[\mathbf M_V]=p^*A.
\]
Every coefficient of $B_V$ is less than one. Apply
Lemma~\ref{prog:stein-representatives} to obtain a Cartier
representative $D_L$ of $L$ on $T_U$. The same birational direct-image
argument applies to the coherent rank-one reflexive canonical sheaf:
it supplies a canonical form nonzero at a chosen point of the common
smooth isomorphism locus. Use this meromorphic form on $T_U$ and its
pullback to $V$ to choose compatible canonical representatives, and set
\[
 N=t p^*D_L-K_V-B_V.
\]
For every $\rho$-vertical curve, its degree equals that of $\mathbf M_V$.
Thus $N$ is $\rho$-nef.

The relative bigness of $N$ uses the additional fact that
$\rho=f\circ p$ is \emph{bimeromorphic}: both the small contraction
$f$ and the resolution $p$ are bimeromorphic. It does not follow from
relative nefness alone. In fact every real Cartier divisor is relatively
big for this projective bimeromorphic map over the Stein neighbourhood.
Here is the required ample-plus-effective decomposition explicitly.
Write $N=\sum_j r_jN_j$, where $r_j>0$, $\sum_jr_j=1$, and the $N_j$
are rational Cartier divisors in its finite Cartier span. Choose a
$\rho$-ample line bundle and a Cartier representative
$P$ for it using Lemma~\ref{prog:stein-representatives}.
If $q_jN_j$ is integral, apply the birational case of that lemma to
$\mathcal O_V(q_jN_j-P)$. Its nonzero holomorphic section has an
effective Cartier zero divisor $E_j$, giving the actual relation
$q_jN_j\sim P+E_j$. This use depends on $\rho$ being bimeromorphic;
the section is allowed to vanish along exceptional fibres. Therefore
\[
 N\sim_{\R}
 \underbrace{\left(\sum_j\frac{r_j}{q_j}\right)P}_{H}
 +\underbrace{\sum_j\frac{r_j}{q_j}E_j}_{E}.
\]
Here $H$ is $\rho$-ample and $E\geq0$, proving $\rho$-bigness.
For a map with positive-dimensional general fibres the direct image
used above can be zero; no such bigness claim is made for that setting.
Make the finitely many section choices on a slightly larger Stein
neighbourhood, and now shrink $U$ so its closure is compact in that
neighbourhood. Properness makes the inverse image of this closure
compact. A log resolution of the finitely many resulting divisors
therefore gives a uniform positive bound for the following choice.
Choose $\epsilon>0$ small enough that $(V,B_V+\epsilon E)$ is sub-klt.
The divisor $(1-\epsilon)N+\epsilon H$ is relatively ample. Express it
as a positive combination of rational ample divisors and take general
members of sufficiently high multiples. Relative Bertini
\cite[Theorem~2.20]{DHP}, applied on a log resolution of $B_V+E$,
gives an effective $\Theta\sim_{\R}N$ such that $(V,B_V+\Theta)$ is
sub-klt. The high multiples make the added coefficients arbitrarily
small. All choices are made near the compact fibre; shrink $U$ once
to retain them.

Put $\Delta=p_*(B_V+\Theta)\geq0$. Write the actual linear equivalence
as $\Theta=N+\sum_j a_j\operatorname{div}_V(g_j)$.
The bimeromorphic map $p$ identifies meromorphic function fields, so
$g_j=p^*h_j$. Pushing down gives
$K_{T_U}+\Delta=tD_L+\sum_j a_j\operatorname{div}_{T_U}(h_j)$;
in particular this divisor is real Cartier. Pulling back this same
identity gives
\[
 K_{T_U}+\Delta\sim_{\R}tD_L,
 \qquad p^*(K_{T_U}+\Delta)=K_V+B_V+\Theta.
\]
Thus $(T_U,\Delta)$ is klt. Any finitely many base line bundles in a
relative linear equivalence can first be trivialized by shrinking $U$.
Record the resulting equality as
\begin{equation}\label{prog:detector-rationalization}
 K_{T_U}+\sum_i d_iD_i=tD_L+\sum_j b_j\operatorname{div}(g_j),
 \qquad d_i>0,
\end{equation}
using the positive support of $\Delta$. After a further relatively
compact shrinking, the displayed divisors have finitely many prime
components: their locally finite supports meet a compact inverse image.
Equality of coefficients in \eqref{prog:detector-rationalization} is a
finite rational affine system in $(d_i,t,b_j)$. The klt condition is
open within this system, as seen on one log resolution; the pullbacks
of its adjoints vary linearly by the displayed identity. A nearby
rational solution therefore gives a rational effective klt boundary
$\Delta_q$ and $t_q\in\Q_{>0}$ with
\begin{equation}\label{prog:detector-local-adjoint}
 K_{T_U}+\Delta_q\sim_{\Q}t_qD_L.
\end{equation}
The individual local primes $D_i$ need not be $\Q$-Cartier: the identity
ensures that the total adjoint is $\Q$-Cartier.

\smallskip\noindent
\emph{One global flip algebra.}
The adjoint in \eqref{prog:detector-local-adjoint} is $f_U$-antiample.
Its ordinary flip and local finite generation follow from
\cite[Theorems~1.14 and~1.18]{FujinoAnalyticBCHM}. Clearing the displayed
linear equivalence identifies a Veronese of its canonical algebra
with a Veronese of $\mathcal R:=\bigoplus_{m\geq0}f_*L^m$ over $U$.
Hence $\mathcal R$ is locally finitely generated by
\cite[Lemma~2.26]{FujinoAnalyticBCHM}. The graded pieces are coherent. Over a connected normal base open
the inverse image is irreducible; products of nonzero sections of line
bundles are nonzero there. Thus its graded section algebra is an
integral domain, without choosing any global meromorphic frame of $L$.
Compactness supplies one sufficiently divisible Veronese
$\mathcal R^{(r)}$ generated in degree one. Its relative Proj $T^+$
restricts to the ordinary flip over every such $U$. It is therefore
normal, locally klt, and small over $Z$, and has the global relatively
ample bundle
$\mathcal O_{T^+}(1)$. On the common big open this bundle is $L^r$.
Reflexivity gives
\[
 (L^+)^{[r]}\simeq\mathcal O_{T^+}(1).
\]
This proves that $L^+$ is a rational line bundle before any assertion
of factoriality. Relative positivity over the compact K\"ahler base
makes $T^+$ compact K\"ahler.

For any global rank-one reflexive sheaf $\mathcal F^+$ on $T^+$,
take its coherent reflexive transform $\mathcal F$ on $T$ using a
common resolution. Some $M=\mathcal F^{[m]}$ is a line bundle.
Choose $c\in\Q$ so that $(M+cL)\cdot C=0$, and clear its denominator
to obtain an integral bundle $J=n(M+cL)$ numerically trivial over $f$.
The local rational klt pairs \eqref{prog:detector-local-adjoint} have
antiample adjoint, so Lemma~\ref{prog:integral-descent} yields
one global line bundle $J_Z$ with $J=f^*J_Z$.
On $T^+$, pull back $J_Z$, undo the rational twist $ncL^+$, and compare
on the common big open. After clearing the already established index
of $L^+$, reflexivity gives an invertible positive reflexive power of
$\mathcal F^+$. This proves global strong factoriality.

\smallskip\noindent
\emph{The original generalized pair.}
The local log-Fano pairs above give $R^if_*\mathcal O_T=0$ for $i>0$
by \cite[Theorem~5.2]{FujinoAnalyticBCHM}. They make $T$ locally klt, hence rational.
Leray for a projective resolution then shows that $Z$ is rational.
Lemma~\ref{prog:bc-descent}(i) therefore gives
$f^*H^{1,1}_{\BC}(Z,\R)=R^\perp$, with injective pullback. Thus $A=t c_1(L)+f^*\eta$ for a unique $\eta$.
Define $A^+=t c_1(L^+)+(f^+)^*\eta$; it is relatively K\"ahler.

To construct its actual discrepancy divisor, let $\mathcal I$ be the
ideal defined by
\[
 \operatorname{im}(f^*f_*L^r\longrightarrow L^r)
   =\mathcal I\otimes L^r.
\]
Principalize $\mathcal I$ on a common resolution $p:W\to T$,
$q:W\to T^+$ carrying $\mathbf M$. If
$\mathcal I\mathcal O_W=\mathcal O_W(-F_L)$, the tautological quotient
gives the actual bundle identity
\[
 q^*\mathcal O_{T^+}(1)=p^*L^r\otimes\mathcal O_W(-F_L).
\]
Hence $E_L=F_L/r$ is effective and exceptional over both sides, and
$[E_L]=p^*c_1(L)-q^*c_1(L^+)$. If $B_W$ is the original structure
boundary, put $B_W^+=B_W-tE_L$. Then
\[
 [K_W+B_W^+]+[\mathbf M_W]=q^*A^+.
\]
Its pushforward is $B^+$, and its discrepancies do not decrease.
The local ordinary comparison in \eqref{prog:detector-local-adjoint}
has discrepancy divisor $t_qE_L$, so its strictness proves the stated
strictness for $tE_L$ as well. Thus the original pair remains gklt.

Finally, write any $\theta\in H^{1,1}_{\BC}(T,\R)$ uniquely as
$\theta=f^*\eta_\theta+s c_1(L)$ and set
$\theta^+=(f^+)^*\eta_\theta+s c_1(L^+)$. Its correction is the actual
divisor $sE_L$. For a rational line $M$, the number $s$ is rational. Applying
Lemma~\ref{prog:integral-descent} to a Cartier multiple of $M-sL$
gives $M=sL+f^*M_Z$ as actual rational lines. Reflexive extension
on the common big open gives $M^+=sL^++(f^+)^*M_Z$.
Consequently the cohomological transport agrees with the reflexive
transform on Chern classes.
\end{proof}


\subsubsection{Real boundaries and their linear transport}

The rational ordinary step constructed below also determines the step
for every sufficiently close real boundary. The relevant comparison is
an identity of actual rational lines over the contraction base.

\begin{proposition}[Ordinary real boundaries]
\label{prog:ordinary-replacement}
Let $T$ be normal compact K\"ahler and globally Weil $\Q$-factorial,
with canonical sheaf a rational line bundle. Let $B\geq0$ be a real
boundary with $(T,B)$ klt. Assume the rational ordinary-step result
of Proposition~\ref{prog:ordinary-step} for $T$ and its subsequent models. Then every
$(K_T+B)$-negative extremal analytic ray has an ordinary step with
projective contractions, compact K\"ahler models, and the expected
opposite relative ample signs. Global Weil $\Q$-factoriality is preserved;
global strong $\Q$-factoriality is preserved when imposed initially.
\end{proposition}
\begin{proof}
Put $D=K_T+B$ and fix a $D$-negative extremal ray $R$. On the finite
positive support of $B$, effectiveness, the klt condition and negativity
on $R$ persist under small coefficient changes. Klt openness is checked
on one log resolution of this support. Choose nearby effective rational
klt boundaries $B_0,\ldots,B_m$ and positive real numbers $u_j$ with
\[
 B=\sum_j u_jB_j,\qquad \sum_j u_j=1,
 \qquad D_j\cdot R<0,\quad D_j=K_T+B_j.
\]
For $B=0$ take just $B_0=0$. Every $D_j$ is a global rational line
bundle. The rational-step theorem for $(T,B_0)$ provides a projective
contraction $f:T\to Z$ onto a compact K\"ahler base, contracting
exactly $R$, with $-D_0$ relatively ample. The base is rational by
\cite[Lemma~8.8]{DHP}, so Lemma~\ref{prog:bc-descent} gives
$f^*H^{1,1}_{\BC}(Z,\R)=R^\perp$; use case~(ii) when $f$ is of
fiber type.
Let $C$ be a rational curve spanning $R$, and set
$a_j=(D_j\cdot C)/(D_0\cdot C)\in\Q_{>0}$. On every projective
fibre, $-D_j$ is numerically equivalent to the positive multiple
$-a_jD_0$. Numerical invariance of ampleness on that fibre, followed
by the relative-ampleness criterion, makes $-D_j$ relatively ample.
The positive combination $-D$ is relatively ample as a real line
bundle. This also handles a fibre-type contraction.

Suppose henceforth that $f$ is birational. A Cartier multiple
$M_j=n_j(D_j-a_jD_0)$ is numerically trivial over $Z$ and
$M_j-D_0$ is relatively ample. Lemma~\ref{prog:integral-descent},
applied to the rational klt pair $(T,B_0)$ gives a rational line bundle
$A_j$ on $Z$ with the actual identity
\begin{equation}\label{prog:real-line-descent}
 D_j=a_jD_0+f^*A_j.
\end{equation}
Here and below an identity of rational line bundles means an isomorphism
after a common integral multiple. The descent lemma is used only for
birational $f$.

For a small $f$, let $f^+:T^+\to Z$ be its rational $D_0$-flip.
The rational theorem supplies a global relatively ample adjoint
$D_0^+$ and preserves the stated factoriality and canonical-sheaf
conditions. The other transformed adjoints $D_j^+$ are therefore
rational line bundles. Transform \eqref{prog:real-line-descent} on the common
big open and extend by reflexivity to obtain
\[
 D_j^+=a_jD_0^++(f^+)^*A_j.
\]
All $D_j^+$ are relatively ample, as is
$D^+=\sum_j u_jD_j^+$. Thus this same small map is the required
ordinary real-boundary flip. The real ordinary discrepancy comparison
proves that $(T^+,B^+)$ is klt and gives the strict increases at
exceptional centres.

For a divisorial contraction, write
$D_0=f^*D_0'+e_0E$ with $e_0>0$. Transforming
\eqref{prog:real-line-descent} on the target and summing yields
\[
 D-f^*D'=\left(\sum_j u_ja_j\right)e_0E.
\]
This coefficient is positive. The same discrepancy comparison proves
the required assertions. The ambient category was already preserved
by the rational step in both cases.
\end{proof}


For later use, let \(D=K_T+B\) drive one of the birational steps just
constructed, and write \(f':T'\to Z\) for its positive morphism
(or the identity for a divisorial contraction). Lemma~\ref{prog:bc-descent}
gives the unique decomposition and forward transform
\begin{equation}
 \theta=f^*\eta+s\{D\},\qquad
 \theta'=(f')^*\eta+s\{D'\}.
 \label{prog:ordinary-forward-transport}
\end{equation}
On a common resolution \(p:W\to T\), \(q:W\to T'\), the ordinary
comparison \(p^*D-q^*D'=F\) is an actual effective real divisor,
exceptional over \(T'\) and over both sides for a flip. Hence
\(p^*\theta-q^*\theta'=s\{F\}\). On Chern classes this is the
reflexive transform, by the line identities in the proof above and
Lemma~\ref{prog:integral-descent}. This construction gives forward
transport without assuming surjectivity onto the next Bott--Chern group.

\subsubsection{Actual real lines and relative numerical spaces}

A second polarization argument will be used for maps whose fibers are
Moishezon. It applies to a real combination of global line bundles,
provided that its degrees on the fibers agree with a relatively Kähler
class.

\begin{lemma}[Polarization by an actual real line bundle]
\label{prog:actual-line-polarization}
Let $f:T\to Z$ be a proper Moishezon map, where $T$ is a compact
K\"ahler space. Suppose that $L\in\operatorname{Pic}(T)\otimes\R$
and a relatively K\"ahler class $\omega$ satisfy
\[
 L\cdot C=\omega\cdot C
 \quad\text{for every compact curve $C$ contracted by $f$.}
\]
Then $L$ is relatively ample and $f$ is projective.
\end{lemma}

\begin{proof}
Let $V$ be an irreducible positive-dimensional subspace of a reduced
fibre. It is Moishezon. Normalize $V$ and take a smooth projective
modification $\pi:\widetilde V\to V^{\mathrm n}$, chosen so that
an effective exceptional divisor $E$ has $-E$ relatively ample.
The pullback of the K\"ahler class from $V$ to its normalization is
K\"ahler. Consequently
$\pi^*\omega-\delta[E]$ is K\"ahler for sufficiently small
$\delta>0$. On the projective manifold $\widetilde V$ the actual
real line bundle $\pi^*L-\delta E$ has the same curve degrees as
that K\"ahler class, and is therefore ample. For example, subtract
a sufficiently small positive multiple of a fixed ample class from
the K\"ahler class; the corresponding real line bundle is nef by
the projective numerical criterion. Also $\pi^*L$ is nef, and the
decomposition
\[
 \pi^*L=(\pi^*L-\delta E)+\delta E
\]
shows that it is big. Thus
\[
 L^{\dim V}\cdot V=(\pi^*L)^{\dim V}>0.
\]
The real Nakai--Moishezon criterion for proper algebraic spaces
\cite[Theorem~1.6]{FujinoMiyamoto2020}, applied to the algebraizations
of the Moishezon fibres, proves that $L$ is ample on every reduced
fibre. The criterion is insensitive to nilpotents and reducible
components. In the coefficient space of the finitely many global
line bundles occurring in $L$, choose, for each fibre, a rational
simplex containing $L$ in its interior whose vertices restrict to
ample rational lines on that fibre. Openness of fibrewise ampleness
for each of these finitely many rational lines gives a base
neighbourhood on which the entire simplex is relatively ample.
Choose finitely many such neighborhoods covering the base and
intersect their coefficient neighborhoods of $L$. Every rational
point of the resulting neighborhood is relatively ample globally.
A rational simplex in this intersection also expresses $L$ as a
positive combination of relatively ample rational lines. Clearing
the denominator of one rational point proves projectivity.
\end{proof}

\begin{remark}
The proof uses only positivity of the top self-intersections of
$L$. Equality of curve degrees with $\omega$ does not assert
equality of their higher intersection numbers or of their
Bott--Chern classes.
\end{remark}




The projective morphism needed for the relative program can be chosen
on a smooth model of the rational quotient. Here the Hodge-theoretic
polarization has an elementary construction, independent of any
birational projectivity-descent assertion.

\begin{lemma}[Smooth Hodge-theoretic polarization]
\label{prog:smooth-polarization}
Let \(g:W\to S\) be a surjective morphism with connected fibers
between smooth compact Kähler manifolds. If
\(g^*:H^0(S,\Omega_S^2)\to H^0(W,\Omega_W^2)\) is an isomorphism,
then \(g\) is projective. This applies when its general fiber is
rationally connected.
\end{lemma}

\begin{proof}
We use the smooth argument of \cite[Theorem~3.1, Step~1]{CH}.
Let \(d=\dim W-\dim S\). Choose a rational class \(u\in H^2(W,\Q)\)
close to a Kähler class. Write its Hodge decomposition as
\(u=b+g^*v\), where \(b\) is Kähler and \(v\) has types \((2,0)\)
and \((0,2)\). Hodge type and the projection formula give
\[
 c=g_*(u^d)=g_*(b^d)>0,\qquad
 g_*(u^{d+1})=g_*(b^{d+1})+(d+1)c\,v.
\]
Since pushforward is rational, the class
\[
 \ell=u-g^*\frac{g_*(u^{d+1})}{(d+1)c}
      =b-g^*\frac{g_*(b^{d+1})}{(d+1)c}
\]
is rational and of type \((1,1)\). Lefschetz's theorem makes it the
Chern class of a rational line. The displayed equality supplies a
metric with positive curvature locally over \(S\), by adding a
potential for the base class. Thus this line is relatively ample.
For rationally connected general fibers, holomorphic two-forms are
pulled back from the smooth base, so the hypothesis holds.
\end{proof}

\begin{lemma}[A rational Hodge correction over a smooth base]
\label{prog:relative-hodge}
Let $f:T\to S$ be an already projective surjective morphism of normal
compact K\"ahler spaces. Assume that $S$ is smooth, $T$ has globally
strongly $\Q$-factorial klt singularities, and, for a projective resolution
$r:W\to T$, the morphism $g=f\circ r$ satisfies
\[
 g^*:H^0(S,\Omega_S^2)\xrightarrow{\ \simeq\ }
                         H^0(W,\Omega_W^2).
\]
Then every $\alpha\in H^{1,1}_{\BC}(T,\R)$ can be written
\[
                         \alpha=c_1(L)+f^*\gamma,
\]
where $L\in\operatorname{Pic}(T)\otimes_{\mathbb Z}\mathbb R$ is a
finite real combination of global line bundles and
$\gamma\in H^{1,1}(S,\mathbb R)$.
\end{lemma}

\begin{proof}
The composition $g$ is projective, since all the spaces are compact.
Choose a $g$-ample line bundle on $W$, with integral Chern class $H$.
Set $d=\dim W-\dim S$ and
\[
                            c=g_*(H^d)>0.
\]
Thus $c$ is the positive degree of $H^d$ on a general fibre. The
cohomological pushforward here is the usual pushforward between the
smooth compact manifolds $W$ and $S$, defined over $\mathbb Q$ and of
Hodge bidegree $(-d,-d)$.

Define a rational linear operator on degree-two cohomology by
\[
 \Pi(\eta)=\eta-g^*\!\left(\frac{g_*(\eta\smile H^d)}{c}\right).
\]
Every class of type $(2,0)$ on $W$ is pulled back from $S$, and the
same holds for type $(0,2)$. The projection formula therefore shows
that $\Pi$ kills both types. It preserves type $(1,1)$. Consequently
\[
        \Pi(H^2(W,\mathbb Q))\subset
                         H^2(W,\mathbb Q)\cap H^{1,1}(W).
\]
By the Lefschetz $(1,1)$ theorem the image consists of rational Chern
classes of global line bundles. Applying $\Pi$ to $r^*\alpha$ gives
\[
 r^*\alpha=c_1(L_W)+g^*\gamma,
 \qquad L_W\in\operatorname{Pic}(W)\otimes\mathbb R,
 \qquad
 \gamma=\frac{g_*(r^*\alpha\smile H^d)}{c}\in H^{1,1}(S,\mathbb R).
\]
For example, express $r^*\alpha$ in a rational basis of $H^2(W,\mathbb Q)$
and apply $\Pi$ to each basis vector. This gives the required finite
real combination $L_W$.

Write $L_W=\sum_j a_jL_j$ with actual line bundles $L_j$ on $W$.
For each $j$, let
\[
                   \mathcal F_j=(r_*L_j)^{**}.
\]
It is a rank-one reflexive coherent sheaf. Strong
$\mathbb Q$-factoriality gives an integer $m_j>0$ for which
$M_j=\mathcal F_j^{[m_j]}$ is a line bundle on $T$.
The canonical identification over the isomorphism locus of $r$ gives
\[
             L_j^{\otimes m_j}\simeq r^*M_j\otimes\mathcal O_W(E_j)
\]
for an integral $r$-exceptional divisor $E_j$.
One may obtain this comparison directly from the evaluation map for
$r_*L_j$: the two coherent rank-one sheaves agree away from the
exceptional locus, so their invertible transforms differ by an
exceptional divisor. No global meromorphic frame of $L_j$ or of the
canonical sheaf is required.

Set $L=\sum_j(a_j/m_j)M_j$ and
$E=\sum_j(a_j/m_j)E_j$. The preceding identities yield
\[
                r^*(\alpha-c_1(L)-f^*\gamma)=[E].
\]
The divisor $E$ is $r$-exceptional and is numerically trivial on every
$r$-contracted curve. The exceptional negativity lemma applied in both
signs gives $E=0$. Injectivity of pullback by a resolution, or pushforward
of the equality of currents, now gives
$\alpha=c_1(L)+f^*\gamma$.
\end{proof}


The hypothesis of Lemma~\ref{prog:relative-hodge} persists on each
projective globally strongly \(\Q\)-factorial klt model over this fixed
smooth base: on a common smooth resolution, holomorphic two-forms are
invariant under modifications. Thus the lemma applies separately on
every working model. It supplies real combinations of global lines;
Lemma~\ref{prog:stein-representatives} supplies their Cartier-divisor
representatives on sufficiently small Stein base neighborhoods.

\begin{lemma}[The projective relative cone in Bott--Chern cohomology]
\label{prog:projective-cone}
Let \(f:T\to S\) be projective between normal compact Kähler spaces
with rational singularities, and assume
\[
 H^{1,1}_{\BC}(T,\R)
   =c_1\bigl(\Pic(T)\otimes\R\bigr)
          +f^*H^{1,1}_{\BC}(S,\R).
\]
Let \(N_1^{\mathrm{line}}(T/S)\) denote the space of real combinations
of \(f\)-contracted curves modulo degrees of global line bundles,
and let \(\overline{\mathrm{NE}}^{\mathrm{line}}(T/S)\) be their
closed effective cone. Their analytic classes give an isomorphism
onto the span of the face
\[
 \mathcal F_f=
 \overline{\mathrm{NA}}(T)\cap(f^*\omega_S)^\perp,
\]
where \(\omega_S\) is any Kähler class on \(S\); this isomorphism
identifies \(\overline{\mathrm{NE}}^{\mathrm{line}}(T/S)\) with
\(\mathcal F_f\). In particular, relative extremal rays in the line
space are extremal rays of the full analytic cone.
\end{lemma}

\begin{proof}
The assumed decomposition makes analytic equivalence and line-degree
equivalence identical on combinations of vertical curves. Their map
into the analytic numerical space is therefore well-defined and
injective. We must show that its closed effective cone is the entire
face, including its classes represented by positive currents.

Fix \(z\in\mathcal F_f\). For every \(\gamma\in
H^{1,1}_{\BC}(S,\R)\), both \(a\omega_S+\gamma\) and
\(a\omega_S-\gamma\) are Kähler for large \(a\). Positivity then
shows that \(z\) annihilates \(f^*\gamma\). Fix a global relatively
ample line \(H\). If a real line \(D\) has nonnegative degree on
every vertical curve, the projective relative numerical criterion
makes \(D+\epsilon H\) relatively ample for every \(\epsilon>0\).
Its Chern class becomes Kähler after adding a sufficiently large
multiple of \(f^*\omega_S\). Therefore
\[
                 (c_1(D)+\epsilon c_1(H))\cdot z\geq0,
 \qquad c_1(D)\cdot z\geq0.
\]
Applying this to both signs of a relatively numerically trivial line
shows that pairing with \(z\) factors through the finite-dimensional
space of relative line degrees. It is nonnegative on its nef cone.
By duality for this projective relative space, it is represented by
an element of \(\overline{\mathrm{NE}}^{\mathrm{line}}(T/S)\).
The decomposition in the statement makes its analytic image equal
to \(z\). Conversely every vertical effective curve lies in
\(\mathcal F_f\), and finite-dimensional injectivity preserves
closedness. This proves the cone equality and the extremality claim.
\end{proof}
